\section{Algoritmet bazë të simulimit probabilitar}
\subsection{Hedhja e monedhës}
\begin{lstlisting}[language=Python,caption={Vlerësimi i probabilitetit dhe gabimit standard.}]
import numpy as np

def monedha(N, p=0.5, seed=2026):
    rng = np.random.default_rng(seed)
    x = rng.random(N) < p
    p_hat = x.mean()
    se = np.sqrt(p_hat*(1-p_hat)/N)
    return p_hat, se, x

for N in [100, 1_000, 10_000, 100_000]:
    p_hat, se, _ = monedha(N)
    print(N, p_hat, se)
\end{lstlisting}
Vlerësimi duhet përsëritur për shumë fara për të parë shpërndarjen e $\hat p$, jo vetëm një trajektore të vetme.

\subsection{Intervali Wilson}
Intervali normal $\hat p\pm z\sqrt{\hat p(1-\hat p)/N}$ sillet dobët për $N$ të vogël ose $\hat p$ pranë 0 e 1. Intervali Wilson është më i qëndrueshëm:
\begin{equation}
 \frac{\hat p+z^2/(2N)\pm z\sqrt{\hat p(1-\hat p)/N+z^2/(4N^2)}}{1+z^2/N}.
\end{equation}

\begin{lstlisting}[language=Python]
def interval_wilson(k, n, z=1.96):
    phat = k/n
    den = 1 + z*z/n
    q = phat + z*z/(2*n)
    d = z*np.sqrt(phat*(1-phat)/n + z*z/(4*n*n))
    return (q-d)/den, (q+d)/den
\end{lstlisting}

\subsection{Përhapja Monte Carlo e pacaktueshmërisë}
\begin{lstlisting}[language=Python,caption={Përhapja e hyrjeve të korreluara.}]
rng = np.random.default_rng(7)
mu = np.array([10.0, 2.0])
Sigma = np.array([[0.25, 0.08], [0.08, 0.04]])
x = rng.multivariate_normal(mu, Sigma, size=200_000)
z = x[:, 0] / x[:, 1]**2
print(np.mean(z), np.std(z), np.quantile(z, [0.025, 0.975]))
\end{lstlisting}
Krahasimi me linearizimin $\nabla f^T\Sigma\nabla f$ tregon kur jolineariteti dhe asimetria bëhen të rëndësishme.

\subsection{Kontrolli i CLT-së}
Për çdo $N$ gjenerohen shumë eksperimente të pavarura dhe studiohet
\begin{equation}
 Z=\frac{\sqrt N(\bar X-\mu)}{\sigma}.
\end{equation}
Histogrami i $Z$ krahasohet me Gauss-in standard. Termi i rekomanduar është \emph{shpërndarje normale} ose \emph{Gauss}, jo përkthime jo të qëndrueshme.

\begin{kujdes}
\emph{Luhatje} përshkruan variacionin e një madhësie të rastit; \emph{gabim standard} përshkruan pacaktueshmërinë e një vlerësuesi. Ato nuk janë sinonime.
\end{kujdes}

Për probabilitetin, statistikën dhe terminologjinë e vlerësimit janë konsultuar materiale universitare shqip \cite{capollari,butka,tasho}; lidhja me simulimin fizik mbështetet te \cite{newman,landau}.
